\begin{abstract}
Turning the equations of a scientific paper into a trustworthy finite-element code is slow and error-prone. Recent systems let a large language model (LLM) write the solver directly; they are reported to produce code that runs but solves the wrong physics, and production C++ libraries such as \dealii{} are among the hardest targets. We present \femagent{}, a framework in which the language model \emph{never writes solver code}. The model transcribes the paper into a typed, physics-aware PDE intermediate representation (fields, tensor ranks, SI dimensions, boundary and initial conditions, constitutive relations, invariants, provenance), which is strictly separated from the numerical problem data and from discretization choices. A deterministic layer then validates the representation, derives weak forms, time discretizations and Gateaux linearizations, generates complete \dealii{} programs (steady and transient; scalar, vector-valued, nonlinear and coupled; with adaptive meshes, Newton solvers and block preconditioning), and \emph{generates the verification}: manufactured solutions, tangent consistency, declared invariants, structural tests, static well-posedness checks, and mutation testing of the verification itself. We evaluate the paper-to-representation step with locally served open-weight models. From a saved web page of a NAFEMS-derived heat-transfer benchmark, a 27B-parameter model yields a verified solver whose value at the reference point agrees with the published one to $2\cdot10^{-4}$ on converged meshes. In a 60-run benchmark (five LLMs, six documents, HTML and PDF text), 55\% of the extractions passed the deterministic checks and 32\% were equivalent to the ground truth, decided by evaluating the extracted mathematics rather than by string matching; of the 14 extractions that passed the checks but were wrong, 13 were stopped by later layers and one produced a solution of a different problem. Extending the prompt contract to the full vocabulary of the representation raised the number of correct extractions from 8 to 14 of 24, and a toy direct-generation baseline did not produce a compiling \dealii{} program. We discuss the limits of the present prototype---a transparent reference backend on rectangular domains, no matrix-free or MPI backend yet, a small and partly synthetic corpus---and the planned extension to matrix-free HPC operators cross-checked against the reference backend.

\medskip\noindent\textbf{Keywords:} finite element method, \dealii{}, large language models, scientific software engineering, intermediate representation, verification, manufactured solutions, mutation testing.
\end{abstract}
